\chapter{Search, Deep Research und Quellenarbeit}

Aktuelle Recherche ist mehr als das Sammeln von Links. Sie beginnt mit einer präzisen Frage, wählt passende Quellen und endet mit Aussagen, deren Herkunft und Unsicherheit nachvollziehbar sind. Je nach Aufgabe reicht eine schnelle Websuche oder es ist eine mehrstufige Tiefenrecherche erforderlich.

\section{Drei Modi der Informationsbeschaffung}

\begin{tabularx}{\textwidth}{@{}lXXX@{}}\toprule
&\textbf{Chatwissen}&\textbf{Websuche}&\textbf{Deep Research}\\\midrule
Tempo&hoch&hoch&mittel\\
Tiefe&begrenzt&punktuell&systematisch\\
Quellen&nicht zwingend&direkte Links&dokumentierter Bericht\\
Einsatz&stabile Grundlagen&aktueller Fakt&komplexe Synthese\\\bottomrule
\end{tabularx}

Chatwissen eignet sich für zeitstabile Erklärungen und erste Orientierung. Websuche ist sinnvoll, sobald sich Informationen verändert haben könnten: Preise, Gesetze, Produktmerkmale, Termine, Rollen oder Softwareversionen. Deep Research eignet sich für Fragen, bei denen viele Quellen gegeneinander abgewogen, Zeiträume eingegrenzt oder widersprüchliche Perspektiven synthetisiert werden müssen.

\section{Eine gute Forschungsfrage formulieren}

Eine offene Themenangabe wie ``KI im Unterricht'' führt leicht zu einer oberflächlichen Sammlung. Eine gute Forschungsfrage enthält Untersuchungsgegenstand, Zielgruppe, Zeitraum und Entscheidungskontext:

\begin{lstlisting}
Welche drei organisatorischen Maßnahmen reduzieren 2026 die größten
Datenschutz- und Qualitätsrisiken beim Einsatz generativer KI an deutschen
Berufsschulen, und welche belastbaren Belege gibt es für ihre Wirksamkeit?
\end{lstlisting}

\section{Quellenhierarchie und Belegqualität}

Primärquellen sind Originaldokumente: Gesetze, Herstellerdokumentation, Forschungsarbeiten, amtliche Statistiken oder direkte Bekanntmachungen. Sekundärquellen ordnen diese ein. Gute Recherche verwendet Sekundärquellen zur Orientierung, stützt zentrale Tatsachenbehauptungen aber möglichst auf Primärquellen.

Für jeden wichtigen Beleg sind vier Fragen hilfreich: Wer ist Urheber? Wann wurde die Information veröffentlicht oder aktualisiert? Welche Methode oder Datengrundlage liegt zugrunde? Belegt die Fundstelle genau die formulierte Aussage oder nur einen ähnlichen Sachverhalt?

\section{Fakten, Interpretation und Unsicherheit trennen}

Eine Quelle kann einen Messwert liefern; die Bedeutung dieses Werts ist bereits Interpretation. Deshalb sollte ein Bericht klar unterscheiden:

\begin{itemize}
\item \textbf{Beobachtung:} direkt aus einer Quelle entnommene Information,
\item \textbf{Einordnung:} Schlussfolgerung aus mehreren Beobachtungen,
\item \textbf{Annahme:} notwendige, aber nicht belegte Arbeitsgrundlage,
\item \textbf{Unsicherheit:} fehlende, widersprüchliche oder methodisch schwache Evidenz.
\end{itemize}

\begin{praxis}
Formuliere eine Forschungsfrage und lege Zeitraum, bevorzugte Primärquellen, Ausschlusskriterien und Ausgabeformat fest. Prüfe anschließend drei zentrale Aussagen direkt an den verlinkten Quellen. Notiere bei jeder Aussage, ob sie wörtlich belegt, abgeleitet oder unsicher ist.
\end{praxis}

